% !TEX spellcheck = en_US



% ~~~~~~~~~~~~~~~~~~~~~~~~~~~~~~~~~~~~~~~ V
% (NC) 2026 Haluk Bingol
% github.com/halukbingol/LaTeX-Templates
%
% Licensees may copy, distribute, display, and perform the work and make derivative works 
% and remixes based on it only for non-commercial purposes. 
% ~~~~~~~~~~~~~~~~~~~~~~~~~~~~~~~~~~~~~~~ A




% =====================================================================
%  Linked Lists -- Examples
%  ---------------------------------------------------------------------
%  Built on hblib/Template/hblibTemplate.tex with the libraries in ../../hblib/.
%  Programs and their real outputs are in codeoutput/:
%     codeoutput/<Name>.java   a program
%     codeoutput/<Name>.txt    what it printed
% =====================================================================

\documentclass[aspectratio=169,10pt,compress]{beamer}

% ~~~~~~~~~~~~~~~~~~~~~~~~~~~~~~~~~~~~~~~~~~~~~~~~~~~~~~~~~~~~~~~~~ V hblib
	\usepackage{../../hblib/hblibPresentation}
	\usepackage{../../hblib/hblibCodeoutput}
	\usepackage{../../hblib/hblibDatastructures}
% ~~~~~~~~~~~~~~~~~~~~~~~~~~~~~~~~~~~~~~~~~~~~~~~~~~~~~~~~~~~~~~~~~ A hblib


\usepackage[T1]{fontenc}
\usepackage{lmodern}
\usepackage{booktabs}

\newcommand{\win}{\textbf{Windows}}
\newcommand{\mac}{\textbf{macOS}}

% long Java lines: wrap instead of running off the frame
\lstset{columns=fixed, breaklines=true, breakatwhitespace=false}

\SetAlCapFnt{\small}
\SetAlCapNameFnt{\small}
\SetKwProg{Fn}{function}{}{end}
\SetKw{KwError}{error}

% ---------------------------------------------------------------------
\title{Linked Lists}
\subtitle{Real Programs}
\author[Bingol]{Ruhan Alpaydin}

\institute[FCIS]{
    Faculty of Computer and Information Sciences\\
    Yeditepe University
}
\date{
	{\tiny \hbTimeStamp}
}

\begin{document}




% ~~~~~~~~~~~~~~~~~~~~~~~~~~~~~~~~~~~~~~ V frm
\begin{frame}[t,fragile]
  \titlepage
\end{frame}
% ~~~~~~~~~~~~~~~~~~~~~~~~~~~~~~~~~~~~~~ A frm




% ~~~~~~~~~~~~~~~~~~~~~~~~~~~~~~~~~~~~~~ V frm
\begin{frame}[t,fragile] \frametitle{Outline}
  \begin{columns}[T]
    \begin{column}{0.48\textwidth}
      \tableofcontents[sections={1-2}, hideallsubsections]
    \end{column}
    \begin{column}{0.48\textwidth}
      \tableofcontents[sections={3-4}, hideallsubsections]
    \end{column}
  \end{columns}
\end{frame}
% ~~~~~~~~~~~~~~~~~~~~~~~~~~~~~~~~~~~~~~ A frm

% ~~~~~~~~~~~~~~~~~~~~~~~~~~~~~~~~~~~~~~ V
\hSection[PPTX]{The Two Exercises from previous years}
% ~~~~~~~~~~~~~~~~~~~~~~~~~~~~~~~~~~~~~~ A




% ~~~~~~~~~~~~~~~~~~~~~~~~~~~~~~~~~~~~~~ V frm
\begin{frame}[t,fragile] \frametitle{Stutter: duplicate every element}
  \begin{columns}[T]
    \begin{column}{0.54\textwidth}
      \lstcode[basicstyle=\ttfamily\tiny]{codeoutput/Stutter.java}{Java}
    \end{column}
    \begin{column}{0.42\textwidth}
      \lstterminal[basicstyle=\ttfamily\scriptsize]{codeoutput/Stutter.txt}
      \small
      \begin{itemize}
        \item Build a new list with each element twice, then swap it in.
        \item \hCode{[1, 2, 3, 4]} $\to$ \hCode{[1, 1, 2, 2, 3, 3, 4, 4]}.
        \item One pass, $O(n)$ time, $O(n)$ extra space.
      \end{itemize}
    \end{column}
  \end{columns}
\end{frame}
% ~~~~~~~~~~~~~~~~~~~~~~~~~~~~~~~~~~~~~~ A frm




% ~~~~~~~~~~~~~~~~~~~~~~~~~~~~~~~~~~~~~~ V frm
\begin{frame}[t,fragile] \frametitle{Mirror: append the reverse of itself}
  \begin{columns}[T]
    \begin{column}{0.54\textwidth}
      \lstcode[basicstyle=\ttfamily\tiny]{codeoutput/Mirror.java}{Java}
    \end{column}
    \begin{column}{0.42\textwidth}
      \lstterminal[basicstyle=\ttfamily\scriptsize]{codeoutput/Mirror.txt}
      \small
      \begin{itemize}
        \item Copy the list, reverse the copy, append it.
        \item \hCode{[1, 2, 3, 4]} $\to$ \hCode{[1, 2, 3, 4, 4, 3, 2, 1]}.
        \item Symmetric input is a palindrome; the deque version needs no copy.
      \end{itemize}
    \end{column}
  \end{columns}
\end{frame}
% ~~~~~~~~~~~~~~~~~~~~~~~~~~~~~~~~~~~~~~ A frm

% ~~~~~~~~~~~~~~~~~~~~~~~~~~~~~~~~~~~~~~ V
\hSection[Classic]{Classic Problems}
% ~~~~~~~~~~~~~~~~~~~~~~~~~~~~~~~~~~~~~~ A




% ~~~~~~~~~~~~~~~~~~~~~~~~~~~~~~~~~~~~~~ V frm
\begin{frame}[t,fragile] \frametitle{Reverse in place}
  \begin{columns}[T]
    \begin{column}{0.54\textwidth}
      \lstcode[basicstyle=\ttfamily\tiny]{codeoutput/ReverseList.java}{Java}
    \end{column}
    \begin{column}{0.42\textwidth}
      \lstterminal[basicstyle=\ttfamily\scriptsize]{codeoutput/ReverseList.txt}
      \small
      \begin{itemize}
        \item Three pointers, one pass: \hCode{prev}, \hCode{cur}, \hCode{next}.
        \item Redirect each \hCode{next} backwards, then advance.
        \item $O(n)$ time, $O(1)$ space, no new nodes.
      \end{itemize}
    \end{column}
  \end{columns}
\end{frame}
% ~~~~~~~~~~~~~~~~~~~~~~~~~~~~~~~~~~~~~~ A frm




% ~~~~~~~~~~~~~~~~~~~~~~~~~~~~~~~~~~~~~~ V frm
\begin{frame}[t,fragile] \frametitle{Palindrome check}
  \begin{columns}[T]
    \begin{column}{0.54\textwidth}
      \lstcode[basicstyle=\ttfamily\tiny]{codeoutput/Palindrome.java}{Java}
    \end{column}
    \begin{column}{0.42\textwidth}
      \lstterminal[basicstyle=\ttfamily\scriptsize]{codeoutput/Palindrome.txt}
      \small
      \begin{itemize}
        \item Two indices, one from each end, meet in the middle.
        \item \hCode{r a c e c a r} $\to$ \hCode{true};
              \hCode{a b c} $\to$ \hCode{false}.
        \item On a linked list without \hCode{get(i)}, the same idea
              needs a stack or a recursion.
      \end{itemize}
    \end{column}
  \end{columns}
\end{frame}
% ~~~~~~~~~~~~~~~~~~~~~~~~~~~~~~~~~~~~~~ A frm




% ~~~~~~~~~~~~~~~~~~~~~~~~~~~~~~~~~~~~~~ V frm
\begin{frame}[t,fragile] \frametitle{Josephus problem}
  \begin{columns}[T]
    \begin{column}{0.54\textwidth}
      \lstcode[basicstyle=\ttfamily\tiny]{codeoutput/Josephus.java}{Java}
    \end{column}
    \begin{column}{0.42\textwidth}
      \lstterminal[basicstyle=\ttfamily\scriptsize]{codeoutput/Josephus.txt}
      \small
      \begin{itemize}
        \item $n$ people in a circle, every $k$-th is removed.
        \item For $n = 7$, $k = 3$, the survivor is $4$.
        \item A circular list is the natural model; the array version
              is $O(n^2)$, the linked version $O(nk)$.
      \end{itemize}
    \end{column}
  \end{columns}
\end{frame}
% ~~~~~~~~~~~~~~~~~~~~~~~~~~~~~~~~~~~~~~ A frm




% ~~~~~~~~~~~~~~~~~~~~~~~~~~~~~~~~~~~~~~ V frm
\begin{frame}[t,fragile] \frametitle{LRU cache with a doubly linked list}
  \begin{columns}[T]
    \begin{column}{0.54\textwidth}
      \lstcode[basicstyle=\ttfamily\tiny]{codeoutput/LruCache.java}{Java}
    \end{column}
    \begin{column}{0.42\textwidth}
      \lstterminal[basicstyle=\ttfamily\scriptsize]{codeoutput/LruCache.txt}
      \small
      \begin{itemize}
        \item Every access moves the entry to the front (most recent).
        \item A \hCode{put} past the capacity evicts the tail
              (least recent).
        \item The JDK ships this as \hCode{LinkedHashMap} with
              \hCode{accessOrder=true} and a \hCode{removeEldestEntry} override.
      \end{itemize}
    \end{column}
  \end{columns}
\end{frame}
% ~~~~~~~~~~~~~~~~~~~~~~~~~~~~~~~~~~~~~~ A frm

% ~~~~~~~~~~~~~~~~~~~~~~~~~~~~~~~~~~~~~~ V
\hSection[Table]{All in One Table}
% ~~~~~~~~~~~~~~~~~~~~~~~~~~~~~~~~~~~~~~ A




% ~~~~~~~~~~~~~~~~~~~~~~~~~~~~~~~~~~~~~~ V frm
\begin{frame}[t,fragile] \frametitle{All the examples in one table}
  \begin{center}\footnotesize
  \begin{tabular}{@{}L{0.14\textwidth}L{0.34\textwidth}L{0.20\textwidth}L{0.18\textwidth}@{}}
    \toprule
    \textbf{Problem} & \textbf{Technique} & \textbf{Time} & \textbf{Space} \\
    \midrule
    Stutter     & rebuild with each element twice           & $O(n)$     & $O(n)$ \\
    Mirror      & copy + reverse + append                    & $O(n)$     & $O(n)$ \\
    ReverseList & three pointers, reverse in place           & $O(n)$     & $O(1)$ \\
    Palindrome  & two indices from the ends                  & $O(n)$     & $O(1)$ \\
    Josephus    & circular list, remove every $k$-th         & $O(nk)$    & $O(n)$ \\
    LruCache    & doubly linked + map, move on access        & $O(1)$     & $O(c)$ \\
    \bottomrule
  \end{tabular}
  \end{center}
  \vspace{0.4em}
  \small
  \hRemark{The two that need no extra memory} are ReverseList and
  Palindrome — both use only a couple of pointers. Every other example
  trades space for simplicity.
\end{frame}
% ~~~~~~~~~~~~~~~~~~~~~~~~~~~~~~~~~~~~~~ A frm




% ~~~~~~~~~~~~~~~~~~~~~~~~~~~~~~~~~~~~~~ V frm
\begin{frame}[t,fragile] \frametitle{Practice}
  \begin{enumerate}\small
    \item Rewrite \hCode{Stutter} so it uses a single node-level
          insertion loop. Does the complexity change?
    \item Rewrite \hCode{Mirror} without copying the list. Which
          extra data structure do you need?
    \item Extend \hCode{ReverseList} to reverse a sub-range
          \hCode{[lo, hi)}.
    \item Extend \hCode{Palindrome} to work on a singly linked list
          without \hCode{get(i)} — first with a stack, then with
          the slow/fast trick from the junior deck.
    \item Extend \hCode{Josephus} to report the \emph{order} of
          removal, not just the survivor.
    \item Implement \hCode{LruCache} from scratch with a
          \hCode{HashMap<K, Node>} and a doubly linked list, no
          \hCode{LinkedHashMap}.
  \end{enumerate}
\end{frame}
% ~~~~~~~~~~~~~~~~~~~~~~~~~~~~~~~~~~~~~~ A frm

% ~~~~~~~~~~~~~~~~~~~~~~~~~~~~~~~~~~~~~~ V
\hSection[Wrap]{Wrap-up}
% ~~~~~~~~~~~~~~~~~~~~~~~~~~~~~~~~~~~~~~ A




% ~~~~~~~~~~~~~~~~~~~~~~~~~~~~~~~~~~~~~~ V frm
\begin{frame}[t,fragile] \frametitle{Big ideas}
  \begin{itemize}
    \item Every example in this deck is a \textbf{node operation}:
          redirect \hCode{next}, \hCode{prev}, or both.
    \item \hRemark{No element is ever moved}; only the links change.
          That is why every operation is $O(1)$ per node touched.
    \item A \textbf{two-pointer walk} (\hCode{prev} + \hCode{cur})
          is the standard way to delete or reverse inside a list.
    \item A \textbf{slow/fast walk} finds the middle or detects a cycle
          in $O(1)$ extra space.
    \item When the problem mentions \emph{recently used} or
          \emph{both ends}, reach for a doubly linked list.
    \item The JDK's \hCode{LinkedList} is a doubly linked list:
          every technique here works on it unchanged.
  \end{itemize}
\end{frame}
% ~~~~~~~~~~~~~~~~~~~~~~~~~~~~~~~~~~~~~~ A frm




% ~~~~~~~~~~~~~~~~~~~~~~~~~~~~~~~~~~~~~~ V frm
\begin{frame}[t,fragile] \frametitle{Check yourself}
  \begin{enumerate}\small
    \item Why does \hCode{ReverseList} need three pointers, not two?
    \item What does \hCode{Stutter} cost in memory if the input list
          has $n$ elements?
    \item Why is the linked-list \hCode{Palindrome} more expensive than
          the array version?
    \item In \hCode{Josephus}, why is the complexity $O(nk)$ rather than $O(n)$?
    \item What data structure does the JDK's \hCode{LinkedHashMap}
          actually use under the hood?
    \item Name the two examples in this deck that use $O(1)$ extra memory.
  \end{enumerate}
\end{frame}
% ~~~~~~~~~~~~~~~~~~~~~~~~~~~~~~~~~~~~~~ A frm




% ~~~~~~~~~~~~~~~~~~~~~~~~~~~~~~~~~~~~~~ V frm
\begin{frame}[t,fragile] \frametitle{References}
  \begin{enumerate}\small
    \item Robert Lafore, \emph{Data Structures \& Algorithms in Java},
          2nd ed., Sams, 2002. (Chapter 5: linked lists.)
    \item Thomas H. Cormen, Charles E. Leiserson, Ronald L. Rivest,
          Clifford Stein, \emph{Introduction to Algorithms},
          4th ed., MIT Press, 2022. (Chapter 10: linked lists.)
    \item Robert Sedgewick and Kevin Wayne, \emph{Algorithms},
          4th ed., Addison-Wesley, 2011. (Section 1.3.)
    \item Flavius Josephus (as retold in Knuth, \emph{The Art of
          Computer Programming}, Vol.~1, Section 1.3.2).
    \item Oracle, \emph{Java SE 21 API Specification}:
          \hCode{java.util.LinkedList},
          \hCode{java.util.LinkedHashMap}.
    \item Oracle, \emph{The Java Tutorials: Trail: Collections}.
  \end{enumerate}
\end{frame}
% ~~~~~~~~~~~~~~~~~~~~~~~~~~~~~~~~~~~~~~ A frm




% ~~~~~~~~~~~~~~~~~~~~~~~~~~~~~~~~~~~~~~ V frm
\begin{frame}[t,fragile]
  \begin{center}
    \vspace{4cm}
    \hRemark{\Huge Questions?}
  \end{center}
\end{frame}
% ~~~~~~~~~~~~~~~~~~~~~~~~~~~~~~~~~~~~~~ A frm




\end{document}